\section{Related Work}\label{sec:related-work}
To address issues of non-executable or non-reproducible research, several solutions have been proposed~(cf.~\cite{siddik2025slr,donat2025r4r}), including tools to gather executable code versions~\cite{DBLP:conf/chi/HeadHBDD19}, recreate the execution order of notebook cells~\cite{DBLP:conf/kbse/WangKLZ20}, provide linting recommendations~\cite{DBLP:journals/ese/PimentelMBF21}, infer the environment of Python code~\cite{DBLP:conf/icse/HortonP19,DBLP:conf/icse/Ye0DW022}, and restore the execution environment of notebooks~\cite{DBLP:conf/icse/WangLZ21}. Yet, neither do these tools address the potential interactions between assumptions, such as the version of an interpreter restricting the available libraries, nor do they provide a systematic approach to document these assumptions in the artifact.
One widely used technique for inferring code constraints is abstract interpretation introduced by \citeauthor{DBLP:conf/popl/CousotC77}~\cite{DBLP:conf/popl/CousotC77}.
It can be used to automatically infer preconditions on the code~\cite{DBLP:conf/vmcai/CousotCFL13} and statically check code contracts~\cite{DBLP:journals/computer/Meyer92,DBLP:conf/foveoos/FahndrichL10}, for example, to infer shapes of input data for data science scripts~\cite{DBLP:journals/corr/abs-2007-10688,urban:hal-04249957}. Additionally, subsequent research by \citeauthor{DBLP:conf/icse/SuboticMS22} presents a general static analysis framework for data science notebooks~\cite{DBLP:conf/icse/SuboticMS22}. Furthermore, \citeauthor{DBLP:conf/pldi/NegriniSU23} propose a static analyzer for data transformations in Jupyter notebooks~\cite{DBLP:conf/pldi/NegriniSU23} and \citeauthor{10.1145/3689609.3689996} introduce a high-level linter to identify code smells in data science scripts~\cite{10.1145/3689609.3689996}. However, none of these approaches focus on identifying and documenting issues of reproducibility, but instead they target the detection of code smells and visualization errors.
